% R11 — Discussion: replace the complete section
\section{Discussion}
\subsection{What the trained representation contributes}
The central finding is that a compact S-JEPA trained on a small collection of pose sequences produces features that support label prediction on held-out source recordings. A linear classifier using these features exceeds the evaluated handcrafted joint-angle RF under both source-weighted and clip-weighted scoring. This makes predictive skeleton learning a useful approach to investigate when large, clinically annotated movement datasets are unavailable.

The training history and downstream predictions provide complementary evidence. The model fits its feature-prediction objective without reducing all inputs to a constant representation, while the frozen-feature classifier measures information associated with the collection labels. A matched untrained encoder would establish how much of that information is gained through pretraining. Mean pose's similar performance also leaves open the contribution of temporal order, as S-JEPA differs from that baseline in visibility inputs, feature dimension and the frames included in its window averages.

\subsection{Scope of the evidence}
The collection was assembled from publicly accessible gait videos, and its Normal, MS and PD labels were not independently clinically verified. Grouping by source video prevents clips from one identified recording from crossing partitions, but the same participant could appear in different recordings. Camera viewpoint, occlusion, pose-estimation errors and other recording characteristics may also contribute to prediction. The visibility-only score of 0.318 shows that detector outputs contain label-associated information; it does not identify whether that information arises from gait or recording conditions.

The RF comparison uses a specific set of lower-limb angle summaries, while S-JEPA uses all 33 landmarks and visibility. It therefore supports a comparison of these implemented pipelines, without establishing an advantage over every handcrafted gait analysis. The small number of training recordings also limits how confidently performance can be extended to new settings. Although the S-JEPA encoder is fitted on these recordings, its MediaPipe pose-estimation front end is pretrained; this is a small-data representation-learning experiment rather than a measurement of the full pipeline's data requirements.

Preprocessing further limits biomechanical interpretation. Centering on the hips and scaling by torso length remove absolute position and apparent body size, so the outputs are not calibrated measurements of walking speed or stride length. Missing detections are interpolated without a maximum gap length. Fifteen short clips are padded by repeating their final frame---five Normal and ten PD clips, with none from the MS collection---which can introduce label-associated temporal patterns. The minimum accepted clip length is eight sampled frames, so a clip need not contain a full gait cycle.

\subsection{Next comparisons}
The most informative next comparison would use the same architecture and readout with a randomly initialized encoder, followed by a supervised model trained on the same sources. Training and evaluating a matched model with shuffled frame order, while holding frames and visibility inputs fixed, would test whether temporal structure contributes to prediction. Comparing anatomical, unbiased regional and motion-based masking would assess the masking choice directly. Evaluation on recordings with verified diagnoses, participant identities and new acquisition settings would then test whether the learned features support reliable behavioral or clinical measurements.
